\chapter{弃牌与跟注}

\section{本章导入}

这是新手最重要、也最反直觉的一章。上一章我们学习了 outs、胜率与底池赔率的基础计算；本章在此基础上，进一步讨论面对下注时如何估算对手范围、选择 bluff-catcher，以及如何用弃牌纪律保护长期 EV。

很多娱乐玩家以为，自己最大的漏洞是不够激进、不敢诈唬、不懂高级打法。事实上，在大多数低级别牌局和娱乐局中，玩家最大的亏损来源往往不是诈唬太少，而是跟注太多。

他们翻前跟注太多，翻牌跟注太多，转牌跟注太多，河牌面对大注仍然跟注太多。只要手里有一对、有听牌、有高牌、有一点希望，就舍不得弃牌。他们常常不是因为理性判断而跟注，而是因为不甘心、好奇、害怕被诈唬、已经投入太多筹码，或者单纯想“看一眼”。

但德州扑克中，跟注不是情绪行为，而是数学行为。每一次跟注，本质上都是在用一定成本购买摊牌或继续实现权益的机会。你必须判断：我需要赢多少比例？我的实际胜率是否足够？对手下注范围中有多少价值牌？有多少诈唬牌？我的牌能击败哪些部分？

如果这些问题回答不清，却只是因为“他可能在诈唬”而跟注，那么长期结果通常会非常糟糕。

弃牌则恰恰相反。很多新手把弃牌看成软弱，认为弃牌就是被欺负、被打跑、错过机会。但成熟牌手知道，弃牌是一种纪律。它不是承认失败，而是承认当前价格、范围和牌力不支持继续。能在错误价格下放弃，能在大底池中放下第二好牌，能克服“不想被诈唬”的心理，是长期盈利的重要能力。

本章的核心输出是“跟注四问”：

\begin{center}
\textbf{我需要多少胜率？\\
对手有哪些价值牌？\\
对手有哪些诈唬牌？\\
我的牌能否击败足够多的下注范围？}
\end{center}

只要你能在面对下注时稳定问这四个问题，就会开始从“好奇型跟注”转向“理性防守”。

\section{核心概念}

\subsection{为什么过度跟注是最大亏损来源}

过度跟注之所以危险，是因为它看起来不像错误。

一次糟糕的诈唬失败，损失很明显；一次错误的大额下注，也容易被记住。但跟注不同。跟注经常披着“合理怀疑”的外衣：我有一对，我有听牌，我赔率好像还行，他也可能在诈唬，我已经投了很多。于是玩家一手又一手地跟，一条街又一条街地跟，最终长期损失大量筹码。

过度跟注会带来几个问题。

第一，你会用太弱的范围进入后续街。翻牌应该弃掉的牌，被你带到了转牌；转牌应该弃掉的牌，被你带到了河牌；河牌应该弃掉的牌，又被你付钱看摊牌。每多跟一条街，错误成本都会增加。

第二，你会让对手的价值下注变得非常赚钱。如果你总是用中等牌力和弱牌支付，对手根本不需要复杂诈唬，只需要不断用强牌下注，就能从你身上拿价值。

第三，你会强化自己的坏习惯。偶尔一次抓到诈唬，会让你记住“我跟对了”；但更多时候你支付给对手价值牌的筹码，却容易被归因于“运气不好”“他刚好有”。这种记忆偏差会让过度跟注持续存在。

第四，你会失去资金和情绪稳定性。河牌错误跟注往往发生在大底池中，一次错误可能抵消很多小底池盈利。连续几次错误跟注后，牌手容易进入追损和上头状态。

娱乐玩家最大的漏洞，常常不是他们完全不会进攻，而是他们无法正确防守。面对下注时，他们跟得太宽、弃得太少、想得太情绪化。要想进步，必须先学会控制跟注频率。

\subsection{跟注不是情绪行为，而是数学行为}

跟注的本质，是用一定成本换取赢下当前底池的机会。它应该是一个数学和范围判断问题，而不是情绪问题。

情绪化跟注常见于以下心理：

\begin{itemize}
\item “我不想被诈唬。”
\item “我已经投了这么多，现在不能弃。”
\item “他不可能每次都有。”
\item “我想看他到底是什么。”
\item “这手牌弃了太弱。”
\item “我今天已经被打走太多次了。”
\end{itemize}

这些想法都不是合格的跟注理由。真正的跟注理由应该是：

\begin{itemize}
\item 当前底池赔率要求我只需要赢一定比例；
\item 我的牌能击败对手下注范围中足够多的组合；
\item 对手有足够多诈唬；
\item 我的牌在防守范围中处于较高位置；
\item 我持有合适阻断牌，减少对手强价值组合；
\item 对手下注尺度和行动线不支持过强价值范围。
\end{itemize}

跟注是一项需要证明的行动。你不能因为“可能赢”就跟注，因为任何牌都有可能赢。你需要判断的是：这个可能性是否足以覆盖跟注成本。

如果你无法说明自己为什么能赢足够多，弃牌通常比跟注更好。

\subsection{底池赔率的基本概念}

底池赔率，是判断跟注是否合理的基础工具。它回答的问题是：

\begin{center}
\textbf{面对当前下注，我需要赢多少比例，跟注才不亏？}
\end{center}

假设底池中已有 100，面对对手下注 50。此时你需要跟注 50，跟注后总底池为 200。你投入 50 去争夺 200 的总底池，因此你需要至少赢：

\[
\frac{50}{200} = 25\%
\]

也就是说，只要你长期能赢超过 25\%，这个跟注在数学上才可能盈利。

再看几个常见例子：

\begin{center}
\begin{tabular}{llll}
\toprule
\textbf{对手下注} & \textbf{原底池} & \textbf{跟注后总底池} & \textbf{所需胜率} \\
\midrule
1/3 底池 & 100 & 166.7 & 约 20\% \\
1/2 底池 & 100 & 200 & 25\% \\
2/3 底池 & 100 & 233.3 & 约 28.6\% \\
1 倍底池 & 100 & 300 & 33.3\% \\
超池 1.5 倍 & 100 & 400 & 37.5\% \\
\bottomrule
\end{tabular}
\end{center}

这个表说明，下注越大，你跟注需要的胜率越高。面对小注，你可以用更宽范围继续；面对大注，你必须显著收紧跟注范围。

但是，底池赔率只是第一步。知道自己需要 25\% 胜率，并不代表你真的有 25\% 胜率。你还需要估算对手下注范围中价值牌和诈唬牌的比例。

\subsection{需要胜率与实际胜率}

跟注决策中，必须区分两个概念：需要胜率和实际胜率。

需要胜率由底池赔率决定。面对 1/2 底池下注，你大约需要 25\% 胜率；面对底池大小下注，你需要约 33.3\% 胜率。

实际胜率则取决于你的牌面对手下注范围的表现。也就是说，你的牌能击败对手多少价值牌和诈唬牌。

很多新手的错误在于，只看到底池很大或下注不算太大，就觉得“赔率不错，可以跟”。但他们没有认真判断自己的实际胜率是否达到要求。

例如河牌底池 100，对手下注 50，你需要约 25\% 胜率。看起来门槛不高。但如果对手下注范围中 90\% 是价值牌，10\% 是诈唬，而你的牌只能击败诈唬，那么你的实际胜率只有 10\%，跟注就是亏损。

反过来，如果对手有大量错过听牌，价值牌组合并不多，你只需要 25\% 胜率，而你的牌能击败 35\% 的下注范围，那么跟注就是合理的。

所以，跟注不能只看价格，也不能只看牌力。必须把价格和范围结合起来：

\begin{center}
\textbf{需要胜率由下注尺度决定；实际胜率由对手范围决定。}
\end{center}

只有实际胜率大于需要胜率，跟注才有意义。

\subsection{面对下注时如何估算对手范围}

面对下注时，新手最常问的是：“他有没有？”但这个问题太模糊。正确问题应该是：

\begin{center}
\textbf{他用哪些价值牌这样下注？他用哪些诈唬牌这样下注？}
\end{center}

估算对手范围，可以按以下步骤进行。

第一，回顾翻前行动。对手是前位加注者、大盲防守者、冷跟注者，还是 3-bet 方？不同翻前行动决定了他的初始范围。

第二，回顾每条街行动。他翻牌下注还是过牌？转牌是否继续？河牌突然下注还是连续开火？每一次行动都会过滤掉一部分牌。

第三，结合牌面变化。哪些听牌完成了？哪些听牌错过了？哪些转牌或河牌强化了对手范围？哪些牌对他的故事不自然？

第四，列出价值牌。包括强顶对、两对、暗三条、顺子、同花、葫芦等。要尽量具体，而不是笼统地说“他可能有强牌”。

第五，列出诈唬牌。包括错过的同花听牌、错过的顺子听牌、没有摊牌价值的浮动牌、带阻断牌的空气牌等。关键是判断这些牌是否真的会用当前尺度下注。

第六，结合对手类型。被动玩家的下注通常偏价值；激进玩家的诈唬更多；跟注站很少在河牌主动 bluff；紧弱玩家大注通常很强。

对手范围不是凭空想象，而是由翻前范围、牌面结构、行动线和对手类型共同决定。

\subsection{哪些牌适合 bluff-catch}

Bluff-catch，即抓诈唬，指的是一手牌通常打不过对手的价值下注，但可以击败对手诈唬。抓诈唬牌的价值不在于它能打赢强牌，而在于对手诈唬足够多时，它可以通过跟注盈利。

适合 bluff-catch 的牌通常具备以下特点：

\begin{itemize}
\item 有一定摊牌价值，可以击败对手空气牌；
\item 不足以主动价值下注；
\item 面对对手价值牌大多落后；
\item 在自己的防守范围中不算太弱；
\item 持有合适阻断牌，减少对手强价值组合；
\item 不阻断对手可能的诈唬组合。
\end{itemize}

例如河牌牌面完成同花，对手大注。你持有一张关键花色 A，但没有成同花。这个 A 可能阻断对手坚果同花组合，使你的抓诈唬价值上升。
又如对手可能错过同花听牌，你手中没有该花色牌，反而不阻断他的错过听牌诈唬组合，这有时也会让跟注更合理。

但新手不要过度迷信 blocker。阻断牌只是辅助因素。最关键的仍然是：对手是否真的有足够诈唬。

很多娱乐玩家河牌大注极少 bluff。面对这类对手，再漂亮的 bluff-catcher 也应该谨慎跟注。

\subsection{哪些牌看似强但应该弃}

新手最难接受的是：有些看似强的牌，在特定行动线和牌面下应该弃。

例如顶对。顶对在翻牌可能很强，但如果对手连续三条街大注，牌面又完成了顺子或同花，你的顶对可能只剩抓诈唬价值。它看起来是顶对，实际上已经无法击败对手大多数价值牌。

例如超对。\hand{A♠ A♣} 在翻前极强，但在 \hand{9♠ 8♣ 7♦ 6♥ T♠} 这样的公共牌上，面对对手持续强力进攻，可能只是非常脆弱的一对。

例如两对。两对在干燥牌面上很强，但在四连顺或三同花牌面上，面对大注可能不再适合支付。

例如小同花。你完成了同花，但如果牌面允许更高同花，对手又表现出极强行动，小同花可能会被更大同花支配。

是否弃牌，不取决于牌名看起来强不强，而取决于它能击败对手当前下注范围的多少部分。

新手需要接受一个事实：

\begin{center}
\textbf{一手牌的绝对牌力很强，不代表它在当前局面中仍然强。}
\end{center}

很多大亏损，正是来自不愿意放下“看起来很强但实际上只是在抓诈唬”的牌。

\subsection{河牌跟注的苛刻条件}

河牌跟注是最苛刻的跟注，因为河牌没有未来权益。你不能靠后续补牌改善，不能再实现听牌，不能寄希望于下一张牌。你跟注后，结果立即揭晓。

因此，河牌跟注必须满足更严格的条件。

第一，你必须知道自己需要多少胜率。下注越大，需要胜率越高。

第二，你必须列出对手的价值牌。如果对手有大量合理价值牌，而你几乎都打不过，那么跟注门槛很高。

第三，你必须列出对手的诈唬牌。如果你找不到足够自然的错过听牌或空气牌，那就不要幻想对手有很多 bluff。

第四，你必须判断对手是否真的会诈唬。很多被动娱乐玩家河牌大注极少诈唬。面对他们，应该过度弃牌，而不是过度跟注。

第五，你的牌必须在防守范围中足够靠前。如果你用所有 bluff-catcher 都跟注，就一定跟太多。你应该选择更好的 bluff-catcher，而不是随便用任何一对支付。

河牌跟注前，最危险的句子是：“我想看一眼。”

看一眼不是免费的。河牌看一眼，往往要付出一大注。你必须确保这次“看一眼”有足够数学依据。

\subsection{如何克服“不想被诈唬”的心理}

“不想被诈唬”是新手过度跟注的核心心理之一。

很多玩家宁愿付钱看错，也不愿弃牌后怀疑自己被 bluff。他们觉得被诈唬是一种羞辱，觉得弃牌就是软弱，觉得对手可能会因此继续欺负自己。于是他们用跟注来维护自尊，而不是维护 EV。

成熟牌手必须理解：被诈唬不是失败，错误跟注才是失败。

如果对手用诈唬赢了你一手小底池，但你长期在不合适的价格下弃牌，你仍然可能是盈利的。反过来，如果你为了不被诈唬而在河牌频繁跟注，你会给对手的价值牌支付大量筹码。

克服这种心理，可以从三个原则开始。

第一，允许自己被诈唬。没有人能抓住所有 bluff。只要你的整体防守频率合理，被个别诈唬打走并不可怕。

第二，关注长期结果，而不是单手自尊。你不是为了证明自己聪明而跟注，而是为了做长期盈利决策。

第三，用范围和价格替代情绪。每次想“我不信”时，改问：“我需要多少胜率？他的诈唬够不够？”

弃牌不是被欺负。弃牌是拒绝在不利价格下支付。

\section{新手常见误区}

\subsection{误区一：已经投了钱，所以不能弃}

这是典型的沉没成本错误。很多新手会想：“我翻牌跟了，转牌也跟了，现在河牌弃掉太可惜。”但已经投入底池的筹码不再属于你。当前决策只应该考虑：现在跟注是否盈利。

如果河牌下注让你需要 33\% 胜率，而你判断自己只有 10\% 胜率，那么即使你前面已经投入很多，也应该弃牌。继续跟注不会把前面的筹码救回来，只会增加新的亏损。

\subsection{误区二：只要对手可能诈唬，就跟注}

对手当然可能诈唬。问题不是“有没有可能”，而是“够不够多”。

如果你面对底池下注，需要约 33\% 胜率。对手下注范围中如果只有 10\% 是诈唬，你跟注就是亏损。即使这一次他刚好在诈唬，也不能证明你的跟注长期正确。

跟注需要的是足够多的诈唬，而不是理论上的可能性。

\subsection{误区三：不看下注尺度}

有些新手面对不同尺度的下注，跟注范围几乎一样。面对 1/3 底池也跟，面对底池下注也跟，面对超池下注还想跟。这是非常严重的错误。

下注越大，你需要胜率越高，对手范围通常也越极化。你必须随着尺度增大而收紧跟注范围。

小注可以用较多边缘牌防守；大注必须用更强的牌和更好的 bluff-catcher 防守。

\subsection{误区四：只看自己的牌名}

“我有顶对。”“我有两对。”“我有同花。”“我有 AA。”这些描述都不够。

跟注时要看的是：我的牌能击败对手下注范围中的哪些牌？如果只能击败诈唬，那它就是 bluff-catcher；如果对手诈唬不足，看起来再强也应该弃。

牌名不能替代范围判断。

\subsection{误区五：跟注为了看信息}

有些玩家河牌跟注的理由是：“我想看看他到底是什么牌。”但信息不是免费的。尤其在河牌，大注看信息的成本非常高。

如果你总是为了信息而跟注，对手只需要用强牌下注，就能从你身上稳定赚钱。你可以通过观察行动线、摊牌记录、对手习惯来获取信息，而不是用错误跟注购买昂贵信息。

\subsection{误区六：不会选择抓诈唬牌}

新手经常随便用一对抓诈唬，却不知道哪些牌更适合 bluff-catch。好的抓诈唬牌通常应当阻断对手价值牌，且不阻断对手诈唬牌。

例如对手可能用错过同花听牌 bluff，如果你手里拿着该花色牌，可能反而减少了他错过同花听牌的组合。相反，如果你没有阻断这些错过听牌，但阻断了他部分强价值牌，跟注质量可能更高。

当然，选择 bluff-catcher 的前提仍然是：对手真的有足够诈唬。

\section{实战思考框架}

本章的核心框架是“跟注四问”：

\begin{center}
\textbf{我需要多少胜率？\\
对手有哪些价值牌？\\
对手有哪些诈唬牌？\\
我的牌能否击败足够多的下注范围？}
\end{center}

\subsection{第一问：我需要多少胜率？}

面对下注，第一步先看价格。你要知道自己跟注需要赢多少比例。

可以先记住几个常见数字：

\begin{center}
\begin{tabular}{ll}
\toprule
\textbf{对手下注尺度} & \textbf{你大约需要的胜率} \\
\midrule
1/3 底池 & 20\% \\
1/2 底池 & 25\% \\
2/3 底池 & 28.6\% \\
1 倍底池 & 33.3\% \\
1.5 倍底池 & 37.5\% \\
\bottomrule
\end{tabular}
\end{center}

这些数字不需要在实战中精确到小数点，但你必须有大致概念。面对越大的下注，你需要越强的理由继续。

\subsection{第二问：对手有哪些价值牌？}

接下来，列出对手可能的价值牌。

不要笼统地说“他可能有强牌”，而要具体：

\begin{itemize}
\item 他能有哪些两对？
\item 他能有哪些暗三条？
\item 他能有哪些顺子？
\item 他能有哪些同花？
\item 他会不会用顶对这样下注？
\item 他前面的行动是否支持这些强牌？
\end{itemize}

列出价值牌后，你要判断自己的牌能击败多少。如果对手价值牌大量存在，而你的牌几乎都打不过，那么跟注必须依赖对手有足够诈唬。

\subsection{第三问：对手有哪些诈唬牌？}

然后，列出对手可能的诈唬牌。

常见诈唬来源包括：

\begin{itemize}
\item 错过的同花听牌；
\item 错过的顺子听牌；
\item 翻牌浮动、转牌获得权益但河牌错过的牌；
\item 持有阻断牌但没有摊牌价值的牌；
\item 由于牌面变化而适合代表强牌的空气牌。
\end{itemize}

你需要判断这些牌是否真实存在于对手范围中。比如一个被动玩家翻牌跟注、转牌跟注、河牌突然大注，他真的会用错过听牌这样打吗？如果不会，你不能强行给他添加大量诈唬。

\subsection{第四问：我的牌能否击败足够多的下注范围？}

最后，把前三个问题结合起来。

如果你需要 25\% 胜率，而对手下注范围中至少有 30\% 是你能击败的牌，跟注可能合理。

如果你需要 33\% 胜率，但对手几乎全是价值牌，你只能击败极少数诈唬，那么应该弃牌。

如果你的牌不仅能击败诈唬，还能击败对手部分薄价值下注，那么跟注质量更高。
如果你的牌只能击败诈唬，那你必须确认对手诈唬足够多。

这就是理性跟注的完整逻辑。

\subsection{弃牌纪律检查表}

当你犹豫是否弃牌时，可以用以下问题提醒自己：

\begin{enumerate}
\item 我是否只是因为不甘心而想跟？
\item 我是否只是因为已经投入筹码而不想弃？
\item 我是否能列出对手足够多的诈唬？
\item 我的牌是否只能击败诈唬？
\item 对手类型是否支持他这样 bluff？
\item 如果我弃牌，即使这次被 bluff，长期是否仍然合理？
\end{enumerate}

如果你发现跟注理由主要来自情绪，而不是范围和价格，弃牌通常是更好的选择。

\section{牌例拆解}

\subsection{牌例一：河牌顶对面对被动玩家大注}

\subsubsection*{牌局背景}

有效筹码 100BB。你在按钮位拿到 \hand{KQo}，open raise 到 2.5BB。大盲是一名偏被动的娱乐玩家，平时很少主动诈唬。

大盲跟注。

翻牌为 \hand{Q♠ 8♣ 2♦ rainbow}。大盲过牌，你下注 1/3 底池，大盲跟注。

转牌为 \hand{5♠}。大盲过牌，你下注 2/3 底池，大盲跟注。

河牌为 \hand{J♠}。大盲突然下注接近底池。

\subsubsection*{新手常见想法}

很多新手会想：

\begin{itemize}
\item “我有顶对好踢脚，不能轻易弃。”
\item “他可能错过听牌在诈唬。”
\item “我已经打了两条街，现在弃太可惜。”
\item “万一他看我害怕河牌 J bluff 呢？”
\end{itemize}

这些想法都需要经过跟注四问检验。

\subsubsection*{理性分析}

第一，我需要多少胜率？面对接近底池下注，你大约需要 33\% 左右的胜率。

第二，对手有哪些价值牌？大盲可能有 \hand{Q♠ J♣} 河牌两对，\hand{8♠ 8♣}、\hand{5♠ 5♣}、\hand{2♠ 2♣} 暗三条，也可能有慢打强牌。被动玩家河牌突然主动大注，通常价值倾向很强。

第三，对手有哪些诈唬牌？翻牌较干燥，转牌也没有明显前门同花听牌。自然错过的听牌并不多。即使有一些 \hand{T♠ 9♣}、\hand{9♠ 7♣} 类牌，也要问：这个偏被动玩家真的会用它们河牌大注吗？答案通常是否定的。

第四，我的牌能否击败足够多下注范围？你的 \hand{K♠ Q♣} 能击败部分诈唬和少数过度下注的弱 Q，但面对对手价值范围大多落后。如果诈唬不足，实际胜率达不到跟注要求。

\subsubsection*{推荐行动}

弃牌。

\subsubsection*{本手牌教训}

顶对好踢脚不是自动跟注牌。面对被动玩家河牌突然大注，必须尊重其价值范围。弃掉顶对不是软弱，而是纪律。

\subsection{牌例二：转牌听牌面对大注}

\subsubsection*{牌局背景}

有效筹码 100BB。你在按钮位拿到 \hand{Q♠ J♠}，open raise 到 2.5BB。大盲跟注。

翻牌为 \hand{A♠ 7♠ 2♦}。大盲过牌，你下注 1/3 底池，大盲跟注。

转牌为 \hand{4♥}。大盲突然下注接近底池。

\subsubsection*{新手常见想法}

新手可能会想：

\begin{itemize}
\item “我有同花听牌，不能弃。”
\item “河牌来黑桃就能赢大底池。”
\item “我翻牌都打了，现在继续很正常。”
\end{itemize}

这手牌的关键，是转牌只剩一张牌，面对大注需要严格计算。

\subsubsection*{理性分析}

第一，我需要多少胜率？面对接近底池下注，你需要约 33\% 胜率。

第二，我的实际权益如何？你有同花听牌，但只剩一张河牌。普通同花听牌完成概率不足以单独支撑面对接近底池的大注。你可能还有 Q、J 高张价值，但面对大盲主动大注，这些高张未必有效。

第三，对手有哪些价值牌？大盲可能有强 A、两对、暗三条，甚至一些转牌改善牌。

第四，对手有哪些诈唬牌？他可能有部分弱听牌或空气牌，但需要判断其是否会主动大注。若对手不是激进玩家，诈唬比例未必足够。

第五，隐含赔率是否足够？即使河牌完成同花，你是否一定能再赢很多？如果大盲看到第三张黑桃后会停止支付，你的隐含赔率就不高。

\subsubsection*{推荐行动}

多数情况下弃牌，除非对手非常激进、下注范围中有大量诈唬，或你有明确隐含赔率。

\subsubsection*{本手牌教训}

有听牌不等于必须跟。转牌听牌面对大注时，只剩一张牌，继续条件比翻牌苛刻得多。

\subsection{牌例三：合格的河牌 bluff-catch}

\subsubsection*{牌局背景}

有效筹码 100BB。你在大盲拿到 \hand{A♥ Q♦}。按钮位 open raise 到 2.5BB，你跟注。

翻牌为 \hand{Q♥ 7♠ 2♥}。你过牌，按钮下注 1/3 底池，你跟注。

转牌为 \hand{4♣}。你过牌，按钮下注 2/3 底池，你跟注。

河牌为 \hand{8♦}。前门红心同花没有完成。你过牌，按钮下注 3/4 底池。

按钮是一名较激进玩家，过去多次在错过听牌时开火。

\subsubsection*{新手常见想法}

新手可能会想：

\begin{itemize}
\item “我有顶对好踢脚，应该跟。”
\item “他可能在 bluff。”
\item “这手牌不能弃。”
\end{itemize}

这次跟注可能合理，但仍然需要完整分析。

\subsubsection*{理性分析}

第一，我需要多少胜率？面对 3/4 底池下注，你需要约 30\% 左右胜率。

第二，对手有哪些价值牌？按钮可能有 \hand{A♠ Q♣}、\hand{K♠ Q♣}、\hand{Q♠ J♣}、\hand{7♠ 7♣}、\hand{2♠ 2♣}、\hand{4♠ 4♣}、少量两对。你的 \hand{A♠ Q♣} 能击败部分较弱 Q，输给暗三条和两对。

第三，对手有哪些诈唬牌？前门红心同花错过。按钮作为激进玩家，可能用错过红心听牌、\hand{J♠ T♣}、\hand{T♠ 9♣}、\hand{6♠ 5♣} 等牌继续 bluff。

第四，阻断牌如何？你持有 \hand{A♥}，阻断了对手部分坚果红心听牌诈唬，也阻断部分强 A 高红心组合。这一点有利有弊。它减少对手一些错过红心 bluff，但也减少对手部分 \hand{A♥ Q♥} 类强价值或强听牌组合。具体要结合对手范围判断。

第五，我的牌能否击败足够多下注范围？你有顶对好踢脚，能击败部分较弱 Q 和诈唬。面对激进玩家，这可能达到跟注要求。

\subsubsection*{推荐行动}

可以跟注。

\subsubsection*{本手牌教训}

抓诈唬不是因为“不想弃”，而是因为对手有足够诈唬、你的牌在防守范围中足够强，并且当前价格支持跟注。

\subsection{牌例四：看似强的同花也可能该弃}

\subsubsection*{牌局背景}

有效筹码 150BB。你在 cutoff 拿到 \hand{8♠ 7♠}，open raise 到 2.5BB。按钮跟注，大盲跟注。

翻牌为 \hand{K♠ T♠ 2♦}。大盲过牌，你下注 1/2 底池，按钮跟注，大盲弃牌。

转牌为 \hand{3♠}，你完成同花。你下注 2/3 底池，按钮跟注。

河牌为 \hand{Q♦}。你下注 1/2 底池，按钮大额加注接近全下。

按钮是一名偏紧被动玩家，很少在河牌大额诈唬。

\subsubsection*{新手常见想法}

很多新手会想：

\begin{itemize}
\item “我已经成同花了，怎么能弃？”
\item “他可能有两对或暗三条乱加注。”
\item “我等了这么久中花，必须看。”
\end{itemize}

但这是一个需要非常谨慎的局面。

\subsubsection*{理性分析}

第一，我需要多少胜率？面对接近全下的大额加注，你需要较高胜率。

第二，对手有哪些价值牌？按钮冷跟注范围中可以有许多更高同花，尤其是 \hand{A♠ X♠}、\hand{Q♠ J♠}、\hand{J♠ 9♠}、\hand{Q♠ 9♠} 等。转牌完成同花后，他跟注；河牌面对你下注大额加注，价值范围极强。

第三，对手有哪些诈唬牌？紧被动玩家河牌大额加注 bluff 非常少。公共牌已经有三张黑桃，他用两对或暗三条转诈唬的可能性较低。

第四，我的牌能击败哪些下注范围？你只有 8 高同花，能击败的价值同花很少，主要只能击败诈唬。但对手诈唬不足。

\subsubsection*{推荐行动}

弃牌。

\subsubsection*{本手牌教训}

完成同花不等于一定能打大底池。低同花面对紧被动玩家河牌大额加注，常常只是 bluff-catcher，甚至是很差的 bluff-catcher。

\subsection{牌例五：不要因为已经投入而河牌跟注}

\subsubsection*{牌局背景}

有效筹码 100BB。你在中位拿到 \hand{A♠ A♣}，open raise 到 3BB。大盲跟注。

翻牌为 \hand{J♠ T♠ 9♦ two-tone}。大盲过牌，你下注 2/3 底池，大盲跟注。

转牌为 \hand{8♠}。大盲过牌，你过牌。

河牌为 \hand{7♠}。大盲下注底池大小。

\subsubsection*{新手常见想法}

新手可能会想：

\begin{itemize}
\item “我有 AA，不能轻易弃。”
\item “我翻牌已经下注了。”
\item “他可能错过同花听牌。”
\item “底池这么大，弃了太可惜。”
\end{itemize}

这里最需要克服的，就是沉没成本和牌名执念。

\subsubsection*{理性分析}

第一，我需要多少胜率？面对底池下注，你需要约 33\% 胜率。

第二，对手有哪些价值牌？公共牌 \hand{J♠ T♣ 9♦ 8♥ 7♠} 极度连接。大盲可以有大量顺子组合，许多两对和暗三条也可能转化为强牌。你的 \hand{A♠ A♣} 在这个牌面上只是单对。

第三，对手有哪些诈唬牌？可能有错过同花听牌，但在如此连接的牌面上，对手有大量自然价值牌。并且大盲底池下注通常表示较强范围。

第四，我的牌能否击败足够多下注范围？你主要只能击败诈唬和少数过度下注的弱牌。若对手不是高频诈唬玩家，跟注不合理。

\subsubsection*{推荐行动}

弃牌。

\subsubsection*{本手牌教训}

\hand{A♠ A♣} 只是翻前最强起手牌，不是所有公共牌上的最强牌。已经投入的筹码不能成为河牌继续支付的理由。

\section{本章总结}

本章讨论的是新手最重要、也最反直觉的能力：弃牌与跟注。

娱乐玩家最大的漏洞，通常不是诈唬太少，而是跟注太多。他们害怕被诈唬，不愿放弃已经投入的筹码，想看对手到底有什么，于是不断用不够强的范围支付。长期来看，这会让对手的价值下注变得极其赚钱。

跟注不是情绪行为，而是数学行为。每一次跟注都必须回答四个问题：

\begin{center}
\textbf{我需要多少胜率？\\
对手有哪些价值牌？\\
对手有哪些诈唬牌？\\
我的牌能否击败足够多的下注范围？}
\end{center}

需要胜率由底池赔率决定，实际胜率由对手范围决定。只有当你的实际胜率大于需要胜率时，跟注才有意义。

河牌跟注尤其苛刻，因为河牌没有未来权益。你不能再靠补牌改善，不能再“看一张”，只能用当前牌力面对对手下注范围。如果你的牌只能击败诈唬，那么你必须确认对手真的有足够多诈唬。

弃牌不是软弱，而是纪律。被诈唬一手牌并不可怕，长期错误跟注才可怕。优秀牌手允许自己被合理诈唬，但不会为了维护自尊，在明显不利的价格下不断支付。

到这里，本书基础部分的核心逻辑已经完成：翻前建立范围，翻牌理解结构，转牌重新评估，河牌做最终判断；再进一步，你需要根据牌力功能决定任务，根据下注目的使用筹码，根据激进条件主动争取底池，并在面对进攻时用理性框架决定跟注或弃牌。

真正的进步，不是学会某一个固定答案，而是建立一套可重复、可复盘、可修正的决策系统。

\section{课后训练}

\subsection*{训练一：跟注四问记录}

下一次打牌后，记录 10 手你面对下注选择跟注的牌。每一手回答：

\begin{enumerate}
\item 我当时需要多少胜率？
\item 对手有哪些价值牌？
\item 对手有哪些诈唬牌？
\item 我的牌能否击败足够多的下注范围？
\item 我的跟注是理性判断，还是情绪驱动？
\end{enumerate}

训练目标是让每一次跟注都有明确理由。

\subsection*{训练二：河牌错误跟注复盘}

找出最近 5 手你河牌跟注后输牌的手牌，逐手分析：

\begin{itemize}
\item 我是否只是在抓诈唬？
\item 对手是否真的有足够诈唬？
\item 我是否忽略了下注尺度？
\item 我是否因为不想被诈唬而跟注？
\item 如果重来一次，我是否应该弃牌？
\end{itemize}

这个训练的目标，是减少河牌高成本错误。

\subsection*{训练三：成功弃牌记录}

记录 5 手你成功弃牌的手牌，哪怕你不知道对手最终是什么牌。

重点写下：

\begin{enumerate}
\item 我为什么选择弃牌？
\item 对手有哪些价值牌？
\item 对手诈唬是否不足？
\item 我的牌是否只是 bluff-catcher？
\item 即使这次被诈唬，长期弃牌是否仍然合理？
\end{enumerate}

这个训练的目标，是强化“弃牌也是正确决策”的正反馈。

\subsection*{训练四：底池赔率练习}

随机选择 20 个下注场景，计算你需要的大致胜率：

\begin{itemize}
\item 对手下注 1/3 底池；
\item 对手下注 1/2 底池；
\item 对手下注 2/3 底池；
\item 对手下注 1 倍底池；
\item 对手下注 1.5 倍底池。
\end{itemize}

不要求精确到小数点，但要形成尺度越大、跟注门槛越高的直觉。

\subsection*{训练五：bluff-catcher 筛选}

找出 10 手你在河牌持有一对或中等牌力面对下注的牌，判断它们是否是好的 bluff-catcher。

逐手回答：

\begin{itemize}
\item 我的牌能否击败对手部分价值下注？
\item 如果不能，它是否只能击败诈唬？
\item 我是否阻断对手强价值牌？
\item 我是否阻断对手诈唬牌？
\item 对手类型是否支持他有足够诈唬？
\end{itemize}

训练目标是学会选择高质量跟注牌，而不是随便用任何一对抓诈唬。

\subsection*{训练六：克服不想被诈唬}

每次你因为“不想被诈唬”而想跟注时，暂停并写下：

\begin{center}
\textbf{被诈唬一手牌不可怕，长期错误跟注才可怕。}
\end{center}

然后回答：

\begin{enumerate}
\item 我是否能列出足够多诈唬组合？
\item 这个对手是否真的会 bluff？
\item 我跟注是为了盈利，还是为了看牌？
\item 如果我弃牌，即使这次被 bluff，长期是否仍然合理？
\end{enumerate}

这个训练的目标，是把跟注从自尊反应，转化为理性决策。
